\documentclass[12pt]{article}

% ======================================================
% Codificación e idioma
% ======================================================

\usepackage[utf8]{inputenc}
\usepackage[T1]{fontenc}
\usepackage{multicol}
\usepackage[spanish,es-nodecimaldot]{babel}

% =====================================================
% Matemáticas
% =====================================================

\usepackage{amsmath,amssymb,amsfonts}

% =====================================================
% Formato y utilidades
% =====================================================

\usepackage[letterpaper,margin=2.2cm]{geometry}
\usepackage{enumitem}
\usepackage{array}
\usepackage{tabularx}
\usepackage{xcolor}
\usepackage{fancyhdr}
\usepackage{lastpage}
\usepackage{hyperref}

% =====================================================
% Hipervínculos
% =====================================================

\hypersetup{
  colorlinks=true,
  linkcolor=blue,
  urlcolor=blue,
  citecolor=blue
}

% =====================================================
% Formato de párrafos
% =====================================================

\setlength{\parindent}{0pt}
\setlength{\parskip}{0.5em}

% =====================================================
% Datos editables del examen
% =====================================================

\newcommand{\institucion}{Tecnológico Nacional de México}
\newcommand{\materia}{Cálculo Diferencial}
\newcommand{\nombreexamen}{Primer Examen Parcial}
\newcommand{\unidad}{Unidad 1: Números Reales}
\newcommand{\docente}{Carlos Ernesto Martínez}
\newcommand{\fechaexamen}{}
\newcommand{\duracion}{60 minutos}
\newcommand{\puntajetotal}{100 puntos}

% =====================================================
% Encabezado y pie de página
% =====================================================

\setlength{\headheight}{15pt}

\pagestyle{fancy}
\fancyhf{}

\fancyhead[L]{\institucion}
\fancyhead[R]{\materia}

\fancyfoot[L]{\docente}
\fancyfoot[C]{Página \thepage\ de \pageref{LastPage}}
\fancyfoot[R]{\nombreexamen}

\renewcommand{\headrulewidth}{0.4pt}
\renewcommand{\footrulewidth}{0.4pt}

% =====================================================
% Contador y ambiente para problemas
% =====================================================

\newcounter{problema}

\newenvironment{problema}[1]
{
  \refstepcounter{problema}
  \par\medskip
  \noindent
  \textbf{Problema \theproblema.}
  \textbf{[#1 puntos]}
  \par\smallskip
}
{
  \par\medskip
}

% =====================================================
% Comando para espacio de respuesta
% =====================================================

\newcommand{\espaciorespuesta}[1]{
  \vspace{#1}
}

% =====================================================
% Documento
% =====================================================

\begin{document}

% =====================================================
% Encabezado principal
% =====================================================

\begin{center}
    {\Large\textbf{\institucion}}\\[0.3em]
    {\large\textbf{\materia}}\\[0.3em]
    {\large\textbf{\nombreexamen}}\\[0.3em]
    {\normalsize\unidad}
\end{center}

\vspace{0.5em}

% =====================================================
% Datos del estudiante
% =====================================================

\begin{tabularx}{\textwidth}{>{\bfseries}l X >{\bfseries}l X}
Nombre: & \hrulefill &
No. de control: & \hrulefill \\[1.2em]

Grupo: & \hrulefill &
Fecha: & \hrulefill \\[1.2em]

Calificación: & \hrulefill &

\end{tabularx}

\vspace{1em}

% =====================================================
% Información general
% =====================================================

% =====================================================
% Instrucciones
% =====================================================

\subsection*{Instrucciones}
Lea cuidadosamente cada problema antes de responder. Muestre de manera clara y ordenada todo el procedimiento. Justifique sus resultados cuando se solicite. Las respuestas sin procedimiento no serán consideradas par la calificaci\'on. Simplifique sus respuestas siempre que sea posible. No se permite el uso de apuntes, sí se permite uso de formulario, mostrar al inicio del examen. Sí se permite uso de calculadora, mostrar al inicio del examen

\medskip


% =====================================================
% REACTIVOS
% =====================================================
\textbf{ Resuelve las siguientes desigualdades, proporcionando el conjunto  solución, el intervalo solución así como la representación gráfica de  la misma}

\begin{multicols}{2}
\begin{problema}{10}
$5x-2<3x+4$
\end{problema}


\begin{problema}{10}
$-2x+1<5x-8$
\end{problema}

\begin{problema}{}
$-\frac{2}{5}x+\frac{1}{3}<-\frac{1}{15}$
\end{problema}

\begin{problema}{10}
$-2x + 1 \geq -5x + 4$
\end{problema}

\begin{problema}{10}
$\frac{1}{4}x - \frac{1}{2} \geq \frac{2}{3}x - \frac{4}{6}$
\end{problema}

\begin{problema}{10}
$-2x + 3 \leq -4x - 1$
\end{problema}

\begin{problema}{10}
$-\frac{1}{2}x + \frac{3}{4} \leq \frac{3}{5}x + \frac{1}{10}$
\end{problema}

\begin{problema}{10}
$-\frac{1}{2}<\frac{3}{4}x+\frac{1}{2}<\frac{1}{2}$
\end{problema}

\begin{problema}{10}
$-\frac{3}{4}<\frac{1}{2}x-\frac{1}{4}<\frac{3}{4}$
\end{problema}

\begin{problema}{10}
$2<-\frac{1}{2}x+\frac{3}{4}\leq 3$
\end{problema}

\begin{problema}{10}
$-\frac{1}{2}\leq \frac{3}{4}x+\frac{1}{2}<\frac{1}{2}$
\end{problema}

\begin{problema}{10}
$-\frac{3}{4}\leq -\frac{1}{2}x+\frac{1}{4}\leq \frac{1}{4}$
\end{problema}

\begin{problema}{10}
$|3x-1|<7$
\end{problema}

\begin{problema}{10}
$|\frac{1}{3}x+\frac{1}{2}|<\frac{5}{6}$
\end{problema}

\begin{problema}{10}
$|2x+1|\leq 7$
\end{problema}

\begin{problema}{10}
$|-\frac{1}{2}x+\frac{3}{4}|\leq 2$
\end{problema}

\begin{problema}{10}
$|3x-1|>4$
\end{problema}


\begin{problema}{10}
$|\frac{1}{2}x-\frac{1}{4}|>\frac{1}{2}$
\end{problema}

\begin{problema}{10}
$|3x-1|\geq 4$
\end{problema}

\begin{problema}{10}
$|\frac{1}{2}x-\frac{1}{4}|\geq\frac{1}{2}$
\end{problema}

\begin{problema}{10}
$|\frac{3}{8}x-\frac{2}{9}| \leq \frac{3}{4}$
\end{problema}

\begin{problema}{10}
$|\frac{1}{6}x-\frac{2}{9}| \leq \frac{1}{4}$
\end{problema}

\begin{problema}{10}
$|{-\frac{1}{5}x+\frac{6}{11}}| > \frac{2}{9}$
\end{problema}

\begin{problema}{10}
$x^2-4x+3<0$
\end{problema}

\begin{problema}{10}
$\frac{1}{2}x^2-\frac{5}{2}x+2\leq0$
\end{problema}

\begin{problema}{20}
$4x^2-12x+9\geq 0$
\end{problema}

\begin{problema}{20}
$x^2+\frac{1}{2}x-\frac{1}{2}\geq 0$
\end{problema}

\begin{problema}{20}
$2x^2-7x+6>0$
\end{problema}

\begin{problema}{20}
$2x^2-5x-3<0$
\end{problema}

\begin{problema}{20}
$-4x^2+16x-12<0$
\end{problema}

\begin{problema}{20}
$2x^2-5x-3<0$
\end{problema}

\begin{problema}{20}
$\frac{1}{3}x^2+\frac{5}{3}x+2\leq0$
\end{problema}
\end{multicols}
\espaciorespuesta{6cm}

% =====================================================
% Tabla de calificación
% =====================================================

\newpage

\subsection*{Registro de calificación}

\begin{center}
\renewcommand{\arraystretch}{1.5}

\begin{tabular}{|c|c|c|}
\hline
\textbf{Problema} & \textbf{Valor} & \textbf{Puntaje obtenido} \\
\hline
1 & 10 & \\
\hline
2 & 15 & \\
\hline
3 & 15 & \\
\hline
4 & 20 & \\
\hline
5 & 20 & \\
\hline
6 & 20 & \\
\hline
\textbf{Total} & \textbf{100} & \\
\hline
\end{tabular}

\end{center}

\vfill

\begin{center}
\textbf{Calificación final:}\qquad
\rule{3cm}{0.4pt}
\end{center}

\end{document}